\documentclass[12pt,a4paper]{article}

\usepackage[T1]{fontenc}
\usepackage[utf8]{inputenc}
\usepackage[top=2.5cm, bottom=2.5cm, left=3cm, right=2.5cm]{geometry}
\usepackage{graphicx}
\usepackage{amsmath}
\usepackage{booktabs}
\usepackage{xcolor}
\usepackage{hyperref}
\usepackage{caption}
\usepackage{enumitem}
\usepackage{titlesec}
\usepackage{float}
\usepackage{parskip}
\usepackage{array}
\usepackage{multirow}
\usepackage{tabularx}

\definecolor{primaryblue}{RGB}{31,78,121}
\definecolor{codeblue}{RGB}{0,114,178}

\hypersetup{
    colorlinks=true,
    linkcolor=primaryblue,
    citecolor=primaryblue,
    urlcolor=codeblue,
    pdftitle={CityPersons Pedestrian Augmentation},
    pdfauthor={BDT-17}
}

\titleformat{\section}{\large\bfseries\color{primaryblue}}{\thesection}{1em}{}[\vspace{-4pt}\rule{\linewidth}{0.6pt}]
\titleformat{\subsection}{\normalsize\bfseries\color{primaryblue!80!black}}{\thesubsection}{1em}{}
\titleformat{\subsubsection}{\normalsize\itshape\color{primaryblue!60!black}}{\thesubsubsection}{1em}{}

\begin{document}

\begin{titlepage}
    \centering
    \vspace*{2cm}
    {\color{primaryblue}\rule{\linewidth}{2pt}}\vspace{0.5cm}

    {\LARGE\bfseries\color{primaryblue}
        CityPersons Pedestrian Augmentation\\[6pt]
        \large using Stable Diffusion 3.5 Medium}

    \vspace{0.4cm}{\color{primaryblue}\rule{\linewidth}{2pt}}\vspace{1.5cm}

    {\large\bfseries Technical Report}

    \vspace{1.2cm}
    \begin{tabular}{ll}
        \textbf{Pipeline}   & Context-Person Composite (V3) \\[4pt]
        \textbf{Backbone}   & Stable Diffusion 3.5 Medium \\[4pt]
        \textbf{Segmenter}  & YOLOv8m-seg \\[4pt]
        \textbf{Hardware}   & Kaggle T4 $\times$2, 448$\times$448 resolution \\[4pt]
        \textbf{Dataset}    & CityPersons (YOLO format, train split) \\
    \end{tabular}

    \vfill
    {\small\color{gray}VinUni AI Thuc Chien -- Track 3: AI Applications \quad|\quad June 2026}
\end{titlepage}

\tableofcontents
\clearpage

\section{Overview}

The V3 Context-Person Composite pipeline generates photorealistic pedestrian augmentations for street scenes by decoupling generation from background preservation. SD3.5 generates candidate pedestrians; YOLOv8m-seg extracts the person mask; only person pixels are composited onto the original frame. This architectural separation ensures background integrity regardless of the model's output.

\section{V3 Architecture}

\subsection{Core Pipeline}
For each sample, V3 executes the following stages:

\begin{enumerate}
    \item \textbf{Smart Placement:} Rule-based bounding-box selection respects road Y-range (0.66--0.92), perspective scale, a $5 \times 5$ slot grid with jitter, and avoids overlaps with existing annotations.

    \item \textbf{Insertion Guide \& Background Neutralisation:} A faint human silhouette ($\alpha=0.44$) is painted at the target. The region is partially blended to neutral grey (strength 0.55), and context is darkened 3\% to maintain texture cues.

    \item \textbf{SD3.5 Img2Img:} Full-resolution frame (448$\times$448) with T5 disabled for VRAM efficiency; model CPU offload and VAE/attention slicing enabled.

    \item \textbf{Person Segmentation:} YOLOv8m-seg runs at confidence threshold 0.12; the detection closest to target with highest composite score is selected. Detection bbox is expanded 2.20$\times$ for spatial search.

    \item \textbf{Mask Validation:} Checks for completeness (height coverage, aspect ratio 1.35--5.2), border contact ($> 6$ px from edge), and minimum area (0.00045 of image).

    \item \textbf{Ghost Detection:} Pixel contrast between source and generated image under mask must exceed 12.0 (0--255 scale). Three checks: early (under mask), mask area, and final (post-paste).

    \item \textbf{Adaptive Retry:} Up to 3 attempts. Rejection reasons trigger specific parameter and prompt adjustments.

    \item \textbf{Appearance Harmonisation:} Six-stage pipeline: local colour transfer (0.72), brightness matching (0.72), contrast matching (0.38, ratio [0.78, 1.16]), saturation matching (ratio [0.72, 1.12]), sensor noise injection (if $>8\%$ cleaner), and sharpness matching (sigma 0.60). Additional colour-temperature pass and edge-halo neutralisation (Gaussian estimate, radius 7.0).

    \item \textbf{Contact Shadow:} Perspective-scaled elliptic shadow blended below feet; opacity scales from 16 (distant) to 46 (near).

    \item \textbf{Output \& Logging:} Final composite, side-by-side PNG, five-panel debug strip (first 24 samples), and manifest CSV with all generation parameters and metadata.
\end{enumerate}

\section{Generation Parameters}

\begin{table}[H]
\centering
\caption{Per-variant hyperparameters.}
\label{tab:gen_params}
\begin{tabular}{lcccc}
\toprule
\textbf{Variant} & \textbf{Strength} & \textbf{Guidance} & \textbf{Steps} & \textbf{Weight} \\
\midrule
Single pedestrian   & 0.72 & 6.8 & 36 & 0.24 \\
Two pedestrians     & 0.74 & 6.9 & 36 & 0.18 \\
Small group (3)     & 0.76 & 7.0 & 38 & 0.16 \\
Occluded pedestrian & 0.74 & 6.8 & 36 & 0.18 \\
Distant pedestrian  & 0.68 & 6.6 & 34 & 0.14 \\
Near pedestrian     & 0.76 & 7.3 & 38 & 0.10 \\
\bottomrule
\end{tabular}
\end{table}

Strength controls deviation from input. Multi-person variants require higher strength (more new content generation). Guidance scale enforces prompt adherence; near-pedestrian uses highest (7.3) because scale and foot contact are hardest to achieve in foreground. Steps balance quality (higher) with runtime (lower); default 36 across variants.

\section{Adaptive Retry Strategy}

\begin{table}[H]
\centering
\caption{Parameter and prompt adjustments per reject reason.}
\label{tab:retry}
\begin{tabular}{p{3.5cm}p{3.5cm}p{4cm}}
\toprule
\textbf{Reason} & \textbf{Parameter change} & \textbf{Prompt adjustment} \\
\midrule
Ghost / low contrast & Strength +0.04, Guidance +0.60 (cap 8.2) & Add ``solid visible person, clear silhouette''; remove ``transparent, faded'' \\[4pt]
Not enough persons & Strength +0.03, Guidance +0.55 (cap 8.6) & Add ``requested count visible, separated silhouettes''; remove ``single person, merged'' \\[4pt]
Too large & Strength $-$0.04, Guidance $-$0.60 (floor 6.0) & Add ``mid-distance pedestrian''; remove ``foreground close-up'' \\[4pt]
Cropped body & Strength $-$0.03, expand context 14\% & Add ``entire body visible, space around body''; remove ``cropped'' \\
\bottomrule
\end{tabular}
\end{table}

Up to 3 retries per sample. Each retry adapts both generation strength/guidance and prompt terms based on the specific rejection reason. Samples exhausting all retries are discarded (no fallback).

\section{Quality Filters}

\subsection{Scale Enforcement}
Two gates:
\begin{itemize}
    \item Early gate (YOLO detection on generated image): height ratio [0.68, 1.32]
    \item Final gate (post-paste): height ratio [0.76, 1.30]
\end{itemize}
Per-variant minimum thresholds: single 0.075, near 0.160, distant 0.065.

\subsection{Ghost-Person Detection}
Three checks prevent low-contrast insertions:
\begin{itemize}
    \item Mean pixel difference under mask $> 12.0$ (0--255 scale)
    \item Segmented mask area $> 0.00045$ of image
    \item Final contrast check on composite $> 0.010$
\end{itemize}

\subsection{Multi-Person Count Validation}
For variants requiring multiple persons, pipeline explicitly verifies expected count was generated. If fewer detected, sample is rejected with reason ``not\_enough\_new\_people'', triggering high-strength adaptive retry.

\subsection{Mask Quality}
Validated for:
\begin{itemize}
    \item Vertical band coverage $> 0.24$ (head, torso, legs all present)
    \item Aspect ratio [1.35, 5.2] (consistent with standing/walking person)
    \item Border contact: masks within 6 px of image edge rejected
    \item Spatial containment: mask outside detection bbox $< 0.22$
\end{itemize}

\section{Smoke Test Results}

\subsection{Configuration}
10 images (multi-person variants only: two-pedestrian and small-group).

\subsection{Outcome}
4 accepted outputs out of 10 attempted (40\% acceptance rate).

\subsection{Reject Reasons}
\begin{itemize}
    \item \textbf{low\_person\_conf} (dominant): After scale failure on attempt 1, adaptive retry reduced strength and guidance, causing further confidence drop on subsequent attempts. \textbf{Linked failure chain:} scale reduction cascades to confidence failure below 0.35 threshold.

    \item \textbf{too\_large\_for\_perspective}: SD3.5 generated persons with height ratio exceeding max (1.32). Adaptive retry correctly reduced strength but lower-energy generation fell below confidence threshold.

    \item \textbf{floating\_or\_bad\_ground}: One rejection for inadequate foot-ground contact (lack of contact shadow or improper positioning).
\end{itemize}

\subsection{Accepted Samples}
\begin{itemize}
    \item bremen\_000209 (two-ped)
    \item aachen\_000123 (two-ped)
    \item bremen\_000109 (two-ped)
    \item bremen\_000186 (small-group)
\end{itemize}

Multi-person small-group had lower acceptance rate than two-pedestrian, consistent with increased difficulty of generating three coherent persons at correct scale.

\section{Technical Findings}

\begin{enumerate}

    \item \textbf{Architectural separation of generation and preservation.} Background preservation cannot be enforced via prompting alone. The only reliable solution is architectural: take background from the original frame; composite only the generated person pixels. This guarantee holds regardless of SD3.5's output.

    \item \textbf{Scale freedom with faint guide outperforms strict constraint.} A loose scale envelope (30\% wider, 16\% taller than target) with faint insertion guide ($\alpha=0.44$) achieves better person quality while still biasing the model toward correct size/location. Strict bounding-box constraints produce ghost silhouettes; loose envelopes permit the model generative freedom.

    \item \textbf{Ghost detection via pixel contrast is non-negotiable.} Without pixel contrast validation, a non-trivial fraction of accepted samples contain near-transparent ``ghost'' persons invisible to downstream detectors. The 12.0 (0--255) threshold is empirically validated.

    \item \textbf{Adaptive retry with reason-specific adjustments outperforms fixed escalation.} A fixed strength increase applied to an over-scaled person makes it even larger on the next attempt. Reason-aware retry (decrease strength for over-scale, expand context for cropped bodies, targeted prompts for ghosts) recovers from failures that fixed escalation cascades on. The smoke test reveals a linked failure chain: scale-retry reduction → confidence failure. Future refinement: widen confidence threshold on retry 2 after scale failure rather than applying full guidance reduction.

    \item \textbf{Multi-person count must be explicitly validated.} A two-pedestrian variant generating one person represents a scenario failure, not a fallback. Explicit count validation prevents silent acceptance of wrong configurations.

\end{enumerate}

\section{Next Steps}

\subsection{Production Run}
\begin{enumerate}
    \item Run 200-image batch across all six variants on train split.
    \item Inspect debug strips for patterns in ghost, scale, and multi-person count rejections.
    \item If \texttt{low\_person\_conf} $> 30\%$ after scale-retry: widen confidence threshold on retry 2 following \texttt{too\_large\_for\_perspective} (accept 0.28 instead of 0.35).
    \item If multi-person count failure $> 40\%$: review insertion guide geometry for small-group case.
    \item Enable post-hoc metrics (SSIM, PSNR, histogram distance, brightness/contrast delta per variant).
\end{enumerate}

\subsection{Quality Gate}
\begin{itemize}
    \item Target: acceptance $\geq 80\%$ and background SSIM $\geq 0.86$ across all variants.
    \item If acceptance $< 60\%$ after parameter sweep: LoRA fine-tuning on CityPersons pedestrian crops.
\end{itemize}

\section{Summary}

V3 resolves background preservation through architectural separation: SD3.5 generates freely, segmentation extracts the person, original background is always used in the final composite. Quality is ensured by scale gates, ghost detection, multi-person count validation, and appearance harmonisation. The smoke test identified a linked failure chain (scale-retry → confidence failure) pointing to specific parameter refinements. A 200-image production run is required to validate against the 80\% / 0.86 SSIM targets.

\begin{thebibliography}{9}

\bibitem{zhang2017citypersons}
S.~Zhang, R.~Benenson, and B.~Schiele,
``CityPersons: A Diverse Dataset for Pedestrian Detection,''
\textit{Proc.\ IEEE CVPR}, 2017.

\bibitem{sd35}
Stability AI, ``Stable Diffusion 3.5 Medium,''
\url{https://stability.ai/news/sd3}, 2024.

\bibitem{yolov8}
G.~Jocher et al., ``Ultralytics YOLOv8,''
\url{https://github.com/ultralytics/ultralytics}, 2023.

\end{thebibliography}

\end{document}
